% =====================================================================
%  Chapter 27 — What a Discovery Would and Would Not Prove
% =====================================================================

\chapter{What a Discovery Would and Would Not Prove}
\label{ch:discovery}

Every discovery announcement in this field's history has been followed by a
period in which the significance of the result was mistaken for the
significance of the conclusion. This chapter is the prophylaxis: what a
genuine detection would establish, what it would not, and what it would take
to convert one into the other.

\section{Four things a detection must state}A claim of dark matter discovery
that does not answer all four of the following is incomplete, and the
incompleteness is where the over-interpretation happens.

The \emph{particle}. Its mass, its nature, and the interaction that
produced the event. A nuclear recoil at 248~keV is consistent with a
scattering event from any particle in a wide mass range, and the mass is
model-dependent through the recoil kinematics.

The \emph{flux}. The number of events expected from the local density,
which requires the velocity distribution of the local halo --- and therefore
a stated $\rho_0$ and a stated velocity profile.

The \emph{significance}, and its dependence on the search. The look-elsewhere
correction, the number of spectra tried, the blinding status, and the
background model in the window.

The \emph{replication}. Whether a second experiment with a different target
and different systematics sees it.

The 248~keV candidate of Chapter~\ref{ch:event248} states the first three
and not the fourth, which is why it is a candidate.

\section{What a particle detection would prove}If a particle were detected with
the four statements above complete, three conclusions would follow, and each
is narrower than the announcement usually implies.

The \emph{existence} of a new massive particle would be established, and
with it the first direct evidence for dark matter as a particle rather than
as a modification of gravity. That is a major result and it would end the
ambiguity of Chapter~\ref{ch:modified}.

The \emph{cross section} at a given mass would be measured, and with it the
local density would become a measurement rather than an assumption, because
the flux and the rate constrain it jointly. This is a quiet but real
by-product, and it would resolve the $0.3$ versus $0.45\ \GeVcmthree$
dispute for the first time.

The \emph{production mechanism} would be constrained by the particle's
properties, and in favourable cases determined. This is where the
origins question of Part~\ref{part:origins} would finally be answerable.

\section{What it would not prove}Three false inferences are worth naming
precisely, because each has been made.

A detected particle \emph{would not} be shown to be all of the dark matter.
The detected particle might account for a fraction of the abundance, with
the remainder in a second component --- the situation in dark matter models
with a dark sector, and the situation that the solar system bound population
would have if it were a distinct species. A claim of ``dark matter
discovery'' that has not measured the abundance fraction is incomplete in the
most consequential way.

A detected particle \emph{would not} settle the small-scale structure
problems. A WIMP detection would leave the cusp-core problem, the missing
satellites and too-big-to-fail exactly where they are, because those are
problems about the density profile of the halo and not about the identity of
the particle.

And a detected particle \emph{would not} explain the origins of the density.
It would constrain the mechanism. The difference between those two verbs is
the difference between Chapter~\ref{ch:origin}'s question and
Chapter~\ref{ch:production}'s, and it is the reason the origins question is
downstream of discovery rather than a reward for it.

\section{What would follow, in order}The consequences of a genuine discovery
would unfold in a predictable order, and it is worth knowing the order in
advance.

In the first months, the community would check the signal against every
known background, and the check would be done by people who did not make the
discovery. This is the only step that is not ambiguous, and it is the step
that should be completed before any interpretation is published.

In the first year, independent experiments would publish their results on
the same region, and the number of those experiments is determined by which
sensors are sensitive to the energy and the target. For a 248~keV nuclear
recoil, that is the deep-underground liquid-xenon and liquid-argon
programmes, and it is why the 248~keV candidate can be settled within the
next-generation cycle rather than the one after it.

Over a few years, the abundance fraction would be measured by the
combination of the detection, the relic density and the production
mechanism, and the field would learn whether it is dealing with one
component or several.

Over a decade, the production mechanism would be pinned down by the
particle's properties and by the cosmological measurements of
Chapter~\ref{ch:cosmoprobes}, and the origins question of
Chapter~\ref{ch:origin} would finally have an answer.

\section{The failure mode to avoid}The failure mode is not the announcement; it is
the interpretation that arrives before the checking is finished. The 3.5~keV
line, the EDGES signal and the 21~cm anomaly were each reported with
interpretations attached, and in each case the interpretation outran the
measurement by months.

The discipline that prevents it is simple and has a cost: a detection is
announced, and the interpretation waits for the replication. That costs the
announcers something and costs the field nothing, because the field has
been wrong four times in five years and has learned that the cost of a
premature interpretation is much larger than the cost of a late one.

\begin{keytakeaways}
\item A complete discovery claim states the particle, the flux, the
search-dependent significance and the replication, and the 248~keV
candidate states the first three only.
\item A genuine detection would establish a particle, constrain the local
density, and constrain the production mechanism --- and would not show that
the particle is all of the dark matter, resolve the small-scale problems, or
explain the origins.
\item The consequences have a predictable order, and the first step is
independent background checking by people who did not make the discovery.
\item The failure mode to avoid is interpretation before replication, and the
discipline against it costs nothing given the field's recent record.
\end{keytakeaways}
